\frfilename{mt2a1.tex}
\versiondate{20.1.13}
\copyrightdate{1996}

\def\chaptername{Lampiran}
\def\sectionname{Teori himpunan}

\def\Enderton{{\smc Enderton 77}}
\def\Halmos{{\smc Halmos 60}}
\def\Henle{{\smc Henle 86}}
\def\Krivine{{\smc Krivine 71}}
\def\Lipschutz{{\smc Lipschutz 64}}
\def\Roitman{{\smc Roitman 90}}

\newsection{2A1}

Khususnya untuk contoh-contoh dalam Bab 21, kita memerlukan beberapa bagian
teori himpunan yang tidak trivial, yang paling baik didekati melalui teori
baku mengenai kardinal dan ordinal; dan di bagian lain jilid ini saya
menggunakan Lema Zorn.
Di sini saya memberikan ikhtisar yang sangat singkat mengenai hasil-hasil
yang terlibat, dengan sebagian besar bukti dihilangkan.
Kebanyakan materi ini seharusnya terdapat dalam setiap pengantar teori
himpunan yang baik.
Rujukan yang saya berikan mengarah ke buku-buku yang kebetulan pernah saya
temui dan dapat saya anjurkan sebagai bahan yang cukup sesuai bagi pemula.

Saya tidak membahas sistem aksioma ataupun landasan logis. Teori himpunan
yang saya gunakan bersifat `naif' dalam arti bahwa saya mengandalkan
pemahaman saya atas pengalaman kolektif seratus tahun terakhir, bukan
suatu upaya deskripsi formal, untuk membedakan argumen yang sah dari argumen
yang tidak aman. Namun, ada beberapa bagian dalam Jilid 5 tempat filsafat
selonggar ini menjadi tidak sesuai, sehingga saya menggunakan argumen yang,
menurut keyakinan saya, dapat diterjemahkan ke dalam teori himpunan
Zermelo--Fraenkel baku tanpa memerlukan gagasan baru.

Walaupun dalam jilid ini saya menggunakan aksioma pilihan tanpa ragu kapan
pun sesuai, saya akan membagi bagian ini menjadi dua bagian, dimulai dengan
gagasan dan hasil yang tidak bergantung pada aksioma pilihan (2A1A--2A1I),
lalu dilanjutkan dengan sisanya (2A1J--2A1P). Saya percaya bahwa bahkan pada
tingkat ini, mempertahankan suatu tingkat pemisahan membantu kita memahami
sifat argumen-argumennya dengan lebih baik.

\leader{2A1A}{Himpunan terurut (a)} Ingat bahwa suatu {\bf himpunan terurut
parsial} adalah suatu himpunan $P$ bersama suatu relasi $\le$ pada $P$
sedemikian sehingga

\inset{jika $p\le q$ dan $q\le r$, maka $p\le r$,

$p\le p$ untuk setiap $p\in P$,

jika $p\le q$ dan $q\le p$, maka $p=q$.}

\noindent Dalam konteks ini, saya akan menuliskan $p\ge q$ untuk menyatakan
$q\le p$, dan $p<q$ atau $q>p$ untuk menyatakan `$p\le q$ dan $p\ne q$'.
$\le$ merupakan suatu {\bf urutan parsial} pada $P$.

\header{2A1A}{\bf (b)} Misalkan $(P,\le)$ suatu himpunan terurut parsial,
dan $A\subseteq P$. Suatu elemen {\bf maksimal} dari $A$ adalah $p\in A$
sedemikian sehingga $p\not<a$ untuk setiap $a\in A$. Perhatikan bahwa $A$
dapat mempunyai lebih dari satu elemen maksimal. Suatu {\bf batas atas} bagi
$A$ adalah $p\in P$ sedemikian sehingga $a\le p$ untuk setiap $a\in A$;
suatu {\bf supremum} atau {\bf batas atas terkecil} adalah batas atas $p$
sedemikian sehingga $p\le q$ untuk setiap batas atas $q$ dari $A$. Paling
banyak hanya ada satu yang demikian\cmmnt{, sebab jika $p$, $p'$ keduanya
merupakan batas atas terkecil, maka $p\le p'$ dan $p'\le p$}. Karena itu kita
dapat dengan aman menuliskan $p=\sup A$ jika $p$ adalah batas atas terkecil
dari $A$.

Demikian pula, suatu elemen {\bf minimal} dari $A$ adalah $p\in A$ sedemikian
sehingga $p\not>a$ untuk setiap $a\in A$; suatu {\bf batas bawah} dari $A$
adalah $p\in P$ sedemikian sehingga $p\le a$ untuk setiap $a\in A$; dan
$\inf A=p$ berarti bahwa

\Centerline{untuk setiap $q\in P$, $p\ge q\iff a\ge q\Forall a\in A$.}

Suatu subhimpunan $A$ dari $P$ disebut {\bf terbatas menurut urutan} jika
mempunyai batas atas dan batas bawah.

Suatu subhimpunan $A$ dari $P$ disebut {\bf terarah ke atas} jika untuk
setiap $p$, $p'\in A$ terdapat $q\in A$ sedemikian sehingga $p\le q$ dan
$p'\le q$\cmmnt{; yaitu, jika setiap subhimpunan berhingga tak kosong dari
$A$ mempunyai batas atas dalam $A$}. Demikian pula, $A\subseteq P$ disebut
{\bf terarah ke bawah} jika untuk setiap $p$, $p'\in A$ terdapat $q\in A$
sedemikian sehingga $q\le p$ dan $q\le p'$\cmmnt{; yaitu, jika setiap
subhimpunan berhingga tak kosong dari $A$ mempunyai batas bawah dalam $A$}.

\cmmnt{Kadang-kadang lebih mudah jika notasi interval tertutup disesuaikan
dengan himpunan terurut parsial sebarang:} $[p,q]$ akan berarti
$\{r:p\le r\le q\}$.

\wheader{2A1A}{0}{0}{0}{24pt}

\header{2A1Ac}{\bf (c)} Suatu {\bf himpunan terurut total} adalah himpunan
terurut parsial $(P,\le)$ sedemikian sehingga

\inset{untuk setiap $p$, $q\in P$, berlaku $p\le q$ atau $q\le p$.}

\noindent Dalam hal ini, $\le$ merupakan urutan {\bf total} atau {\bf linear}
pada $P$.

Dalam setiap himpunan terurut total kita mempunyai suatu {\bf fungsi median}:
untuk $p$, $q$, $r\in P$, tetapkan

$$\eqalign{\med(p,q,r)
&=\max(\min(p,q),\min(p,r),\min(q,r))\cr
&=\min(\max(p,q),\max(p,r),\max(q,r))\dvro{.}{,}\cr}$$

\cmmnt{\noindent sehingga $\med(p,q,r)=q$ jika $p\le q\le r$.}

\header{2A1Ad}{\bf (d)} Suatu {\bf kisi} adalah himpunan terurut parsial
$(P,\le)$ sedemikian sehingga

\inset{untuk setiap $p$, $q\in P$, $p\vee q=\sup\{p,q\}$ dan
$p\wedge q=\inf\{p,q\}$ terdefinisi dalam $P$.}

\header{2A1Ae}{\bf (e)} Suatu {\bf himpunan terurut baik} adalah himpunan
terurut total $(P,\le)$ sedemikian sehingga\cmmnt{ $\inf A$ ada dan termasuk
dalam $A$ untuk setiap himpunan tak kosong $A\subseteq P$; yaitu,} setiap
subhimpunan tak kosong dari $P$ mempunyai elemen terkecil. Dalam hal ini
$\le$ merupakan suatu {\bf pengurutan baik} dari $P$.

\leader{2A1B}{Rekursi Transfinit: Teorema} Misalkan $(P,\le)$ suatu himpunan
terurut baik dan $X$ sebarang kelas. Untuk $p\in P$, tuliskan $L_p$ untuk
himpunan $\{q:q\in P,\,q<p\}$ dan $X^{L_p}$ untuk kelas semua fungsi dari
$L_p$ ke $X$. Misalkan $F:\bigcup_{p\in P}X^{L_p}\to X$ sebarang fungsi.
Maka terdapat tepat satu fungsi $f:P\to X$ sedemikian sehingga
$f(p)=F(f\restr L_p)$ untuk setiap $p\in P$.

\proof{ Terdapat versi-versi hasil ini dalam \Enderton\ (hlm.\ 175) dan
\Halmos\ (\S18). Walaupun demikian, saya menuliskan suatu bukti, sebab
menurut saya kebanyakan buku elementer tentang teori himpunan tidak
memberinya tempat yang semestinya pada awal sekali teori himpunan terurut
baik.

\medskip

{\bf (a)} Misalkan $\Phi$ adalah kelas semua fungsi $\phi$ sedemikian sehingga

\quad ($\alpha$) $\dom\phi$ merupakan subhimpunan $P$, dan
$L_p\subseteq\dom\phi$ untuk setiap $p\in\dom\phi$;

\quad ($\beta$) $\phi(p)\in X$ untuk setiap $p\in\dom\phi$, dan
$\phi(p)=F(\phi\restr L_p)$ untuk setiap $p\in\dom\phi$.

\medskip

{\bf (b)} Jika $\phi$, $\psi\in\Phi$, maka $\phi$ dan $\psi$ sama pada
$\dom\phi\cap\dom\psi$. \Prf\Quer\ Jika tidak, maka
$A=\{q:q\in\dom\phi\cap\dom\psi,\,\phi(q)\ne\psi(q)\}$ tak kosong. Karena
$P$ terurut baik, $A$ mempunyai elemen terkecil, sebutlah $p$. Sekarang
$L_p\subseteq\dom\phi\cap\dom\psi$ dan $L_p\cap A=\emptyset$, sehingga

\Centerline{$\phi(p)=F(\phi\restr L_p)=F(\psi\restr L_p)=\psi(p)$,}

\noindent yang mustahil.\ \Bang\Qed

\medskip

{\bf (c)} Akibatnya, $\Phi$ merupakan suatu himpunan, sebab fungsi
$\phi\mapsto\dom\phi$ adalah fungsi injektif dari $\Phi$ ke $\Cal PP$, dan
inversnya adalah suatu surjeksi dari subhimpunan $\Cal PP$ ke $\Phi$.
Karena itu, tanpa ragu-ragu kita dapat mendefinisikan suatu fungsi $f$ dengan
menuliskan

\Centerline{$\dom f=\bigcup_{\phi\in\Phi}\dom\phi$,
\quad$f(p)=\phi(p)$ kapan pun $\phi\in\Phi$ dan $p\in\dom\phi$.}

\noindent (Jika Anda menganggap suatu fungsi $\phi$ hanya sebagai himpunan
pasangan terurut $\{(p,\phi(p)):p\in\dom\phi\}$, maka $f$ menjadi
$\bigcup\Phi$.) Maka $f\in\Phi$. \Prf\ Tentu saja $f$ adalah suatu fungsi
dari subhimpunan $P$ ke $X$. Jika $p\in\dom f$, maka terdapat $\phi\in\Phi$
sedemikian sehingga $p\in\dom\phi$, dan dalam hal ini

\Centerline{$L_p\subseteq\dom\phi\subseteq\dom f$,
\quad$f(p)=\phi(p)=F(\phi\restr L_p)=F(f\restr L_p)$. \Qed}

\medskip

{\bf (d)} $f$ terdefinisi di seluruh $P$. \Prf\Quer\ Jika tidak,
$P\setminus\dom f$ tak kosong dan mempunyai elemen terkecil, sebutlah $r$.
Sekarang $L_r\subseteq\dom f$. Definisikan suatu fungsi $\psi$ dengan
menyatakan bahwa $\dom\psi=\{r\}\cup\dom f$, $\psi(p)=f(p)$ untuk
$p\in\dom f$ dan $\psi(r)=F(f\restr L_r)$. Maka $\psi\in\Phi$, sebab jika
$p\in\dom\psi$

\inset{{\it atau} $p\in\dom f$, sehingga
$L_p\subseteq\dom f\subseteq\dom\psi$ dan

\Centerline{$\psi(p)=f(p)=F(f\restr L_p)=F(\psi\restr L_p)$}}

\inset{{\it atau} $p=r$, sehingga
$L_p=L_r\subseteq\dom f\subseteq\dom\psi$ dan

\Centerline{$\psi(p)=F(f\restr L_r)=F(\psi\restr L_r)$.}}

\noindent Karena itu $\psi\in\Phi$ dan
$r\in\dom\psi\subseteq\dom f$.\ \Bang\Qed

\medskip

{\bf (e)} Jadi $f:P\to X$ adalah suatu fungsi sedemikian sehingga
$f(p)=F(f\restr L_p)$ untuk setiap $p$. Untuk melihat bahwa $f$ unik,
perhatikan bahwa setiap fungsi jenis ini harus termasuk dalam $\Phi$, sehingga
harus sama dengan $f$ pada domain bersama keduanya, yakni seluruh $P$.
}%end of proof of 2A1B

\cmmnt{\medskip

\noindent{\bf Catatan} Jika Anda telah diajarkan untuk membedakan kata
`himpunan' dan `kelas', Anda akan melihat bahwa teori himpunan naif saya
relatif toleran karena bersedia mengizinkan variabel kelas dalam
teorema-teoremanya.
}

\vleader{48pt}{2A1C}{Ordinal} Suatu {\bf ordinal} (kadang-kadang disebut
`ordinal von Neumann') adalah suatu himpunan $\xi$ sedemikian sehingga

\inset{jika $\eta\in\xi$, maka $\eta$ merupakan suatu himpunan dan
$\eta\not\in\eta$,

jika $\eta\in\zeta\in\xi$, maka $\eta\in\xi$,

dengan menuliskan `$\eta\le\zeta$' untuk menyatakan `$\eta\in\zeta$ atau
$\eta=\zeta$', $(\xi,\le)$ terurut baik\dvro{.}{}}

\prooflet{
\noindent(\Enderton, hlm.\ 191; \Halmos, \S19; \Henle, hlm.\ 27; \Krivine,
hlm.\ 24; \Roitman, 3.2.8.\cmmnt{ Tentu saja banyak teori himpunan tidak
mengizinkan himpunan menjadi anggota dirinya sendiri, dan/atau menganggap
begitu saja bahwa setiap objek pembicaraan merupakan suatu himpunan, tetapi
saya memilih untuk tidak mengambil pendirian mengenai hal-hal semacam itu
secara umum.})
}

\leader{2A1D}{Fakta-fakta dasar mengenai ordinal (a)} Jika $\xi$ adalah
suatu ordinal, maka setiap anggota $\xi$ adalah ordinal. \prooflet{(\Enderton,
hlm.\ 192; \Henle, 6.4;
\Krivine, hlm.\ 14; \Roitman, 3.2.10.)}

\header{2A1Db}{\bf (b)} Jika $\xi$, $\eta$ adalah ordinal, maka berlaku
$\xi\in\eta$ atau $\xi=\eta$ atau $\eta\in\xi$ (dan tidak ada dua dari
ketiganya yang dapat terjadi bersamaan). \prooflet{(\Enderton,
hlm.\ 192; \Henle, 6.4; \Krivine, hlm.\ 14; \Lipschutz, 11.12; \Roitman,
3.2.13.)} Lazimnya\cmmnt{, dalam hal ini,} kita menuliskan $\eta<\xi$ jika
$\eta\in\xi$ dan $\eta\le\xi$ jika $\eta\in\xi$ atau $\eta=\xi$.
Perhatikan bahwa $\eta\le\xi$ jika dan hanya jika $\eta\subseteq\xi$.

\header{2A1Dc}{\bf (c)} Jika $A$ sebarang kelas tak kosong dari ordinal,
maka terdapat $\alpha\in A$ sedemikian sehingga $\alpha\le\xi$ untuk setiap
$\xi\in A$. \prooflet{(\Henle, 6.7; \Krivine, hlm.\ 15.)}

\header{2A1Dd}{\bf (d)} Jika $\xi$ suatu ordinal, maka $\xi\cup\{\xi\}$
juga ordinal; sebutlah ` $\xi+1$'. \cmmnt{Jika $\xi<\eta$, maka
$\xi+1\le\eta$;} $\xi+1$ adalah ordinal terkecil yang lebih besar dari
$\xi$. \prooflet{(\Enderton, hlm.\ 193; \Henle, 6.3;
\Krivine, hlm.\ 15.)} Untuk setiap ordinal $\xi$, terdapat ordinal terbesar
$\eta<\xi$, dan dalam hal ini $\xi=\eta+1$ serta kita menyebut $\xi$ suatu
{\bf ordinal penerus}, atau $\xi=\bigcup\xi$, dan dalam hal ini kita menyebut
$\xi$ suatu {\bf ordinal limit}.

\header{2A1De}{\bf (e)} Beberapa ordinal pertama adalah $0=\emptyset$,
$1=0+1=\{0\}=\{\emptyset\}$,
$2=1+1=\{0,1\}=\{\emptyset,\{\emptyset\}\}$,
$3=2+1=\{0,1,2\}$, $\ldots$. Ordinal tak hingga pertama adalah
$\omega=\{0,1,2,\ldots\}$, yang dapat diidentifikasi dengan $\Bbb N$.

\header{2A1Df}{\bf (f)} Gabungan sebarang himpunan ordinal merupakan suatu
ordinal. \prooflet{(\Enderton, hlm.\ 193; \Henle, 6.8; \Krivine, hlm.\ 15;
\Roitman, 3.2.19.)}

\header{2A1Dg}{\bf (g)} Jika $(P,\le)$ sebarang himpunan terurut baik, maka
terdapat tepat satu ordinal $\xi$ sedemikian sehingga $P$ isomorfik dengan
$\xi$ sebagai himpunan terurut, dan isomorfisme urutan tersebut unik.
\prooflet{(\Enderton, hlm.\ 187--189; \Henle, 6.13; \Halmos, \S20.)}

\leader{2A1E}{Ordinal awal} Suatu {\bf ordinal awal} adalah ordinal $\kappa$
sedemikian sehingga tidak terdapat bijeksi antara $\kappa$ dan sebarang
anggota $\kappa$. \prooflet{(\Enderton, hlm.\ 197; \Halmos, \S25; \Henle,
hlm.\ 34; \Krivine, hlm.\ 24; \Roitman, 5.1.10, hlm.\ 79).}

\leader{2A1F}{Fakta-fakta dasar mengenai ordinal awal (a)} Semua ordinal
berhingga, dan ordinal tak hingga pertama $\omega$, merupakan ordinal awal.

\header{2A1Fb}{\bf (b)} Untuk setiap himpunan terurut baik $P$, terdapat
tepat satu ordinal awal $\kappa$ sedemikian sehingga terdapat bijeksi antara
$P$ dan $\kappa$.

\header{2A1Fc}{\bf (c)} Untuk setiap ordinal $\xi$, terdapat ordinal awal
terkecil yang lebih besar daripada $\xi$. \prooflet{(\Enderton, hlm.\ 195;
\Henle, 7.2.1.)} Jika $\kappa$ suatu ordinal awal, tuliskan $\kappa^+$ untuk
ordinal awal terkecil yang lebih besar daripada $\kappa$. Kita menuliskan
$\omega_1$ untuk $\omega^+$, $\omega_2$ untuk $\omega_1^+$, dan seterusnya.

\header{2A1Fd}{\bf (d)} Untuk setiap ordinal awal $\kappa\ge\omega$ terdapat
bijeksi antara $\kappa\times\kappa$ dan $\kappa$; akibatnya terdapat bijeksi
antara $\kappa$ dan $\kappa^r$ untuk setiap $r\ge 1$.

\leader{2A1G}{Teorema Schr\"oder--Bernstein} Saya\cmmnt{ mengingatkan Anda
akan hasil mendasar berikut: j}ika $X$ dan $Y$ adalah himpunan dan terdapat
injeksi $f:X\to Y$, $g:Y\to X$, maka terdapat bijeksi $h:X\to Y$.
\prooflet{(\Enderton, hlm.\ 147; \Halmos, \S22; \Henle, 7.4;
\Lipschutz, hlm.\ 145; \Roitman, 5.1.2. \cmmnt{Ini juga merupakan kasus
khusus dari 344D dalam Jilid 3.})}

\leader{2A1H}{Subhimpunan terhitung dari $\Cal P\Bbb N$} \cmmnt{Hasil-hasil
berikut akan diperlukan di bawah.

\medskip

}{\bf (a)} Terdapat bijeksi antara $\Cal P\Bbb N$ dan $\Bbb R$.
\cmmnt{(\Enderton, hlm.\ 149; \Lipschutz, hlm.\ 146.)}

\spheader 2A1Hb Andaikan bahwa $X$ sebarang himpunan sedemikian sehingga
terdapat injeksi dari $X$ ke $\Cal P\Bbb N$. Misalkan $\Cal C$ adalah
himpunan subhimpunan terhitung dari $X$. Maka terdapat surjeksi dari
$\Cal P\Bbb N$ ke $\Cal C$. \prooflet{\Prf\ Misalkan
$f:X\to\Cal P\Bbb N$ suatu injeksi. Tetapkan
$f_1(x)=\{0\}\cup\{i+1:i\in f(x)\}$; maka $f_1:X\to\Cal P\Bbb N$ injektif
dan $f_1(x)\ne\emptyset$ untuk setiap $x\in X$. Definisikan
$g:\Cal P\Bbb N\to\Cal PX$ dengan menetapkan

\Centerline{$g(A)=\{x:\exists\,n\in\Bbb N,\,f_1(x)
=\{i:2^n(2i+1)\in A\}\}$}

\noindent untuk setiap $A\subseteq\Bbb N$. Maka $g(A)$ terhitung, sebab kita
mempunyai injeksi

\Centerline{$x\mapsto\min\{n:f_1(x)=\{i:2^n(2i+1)\in A\}\}$}

\noindent dari $g(A)$ ke $\Bbb N$. Jadi $g$ adalah suatu fungsi dari
$\Cal P\Bbb N$ ke $\Cal C$. Untuk melihat bahwa $g$ surjektif, perhatikan
bahwa $\emptyset=g(\emptyset)$, sedangkan jika $C\subseteq X$ terhitung dan
tak kosong, terdapat surjeksi $h:\Bbb N\to C$; sekarang tetapkan

\Centerline{$A=\{2^n(2i+1):n\in\Bbb N,\,i\in f_1(h(n))\}$,}

\noindent dan lihat bahwa $g(A)=C$.\ \Qed}

\spheader 2A1Hc Andaikan lagi bahwa $X$ adalah suatu himpunan sedemikian
sehingga terdapat injeksi dari $X$ ke $\Cal P\Bbb N$, dan tuliskan $H$ untuk
himpunan fungsi $h$ sedemikian sehingga $\dom h$ adalah subhimpunan terhitung
dari $X$ dan $h$ mengambil nilai dalam $\{0,1\}$. Maka terdapat surjeksi dari
$\Cal P\Bbb N$ ke $H$. \prooflet{\Prf\ Misalkan $\Cal C$ adalah himpunan
subhimpunan terhitung dari $X$ dan misalkan $g:\Cal P\Bbb N\to\Cal C$ suatu
surjeksi, seperti dalam (b). Untuk $A\subseteq\Bbb N$, tetapkan

\Centerline{$g_0(A)=g(\{i:2i\in A\})$,\quad$g_1(A)=g(\{i:2i+1\in A\})$,}

\noindent sehingga $g_0(A)$, $g_1(A)$ adalah subhimpunan terhitung dari $X$,
dan $A\mapsto(g_0(A),g_1(A))$ merupakan surjeksi dari $\Cal P\Bbb N$ ke
$\Cal C\times\Cal C$. Misalkan $h_A$ adalah fungsi dengan domain
$g_0(A)\cup g_1(A)$ sedemikian sehingga $h_A(x)=1$ jika $x\in g_1(A)$,
$0$ jika $x\in g_0(A)\setminus g_1(A)$. Maka $A\mapsto h_A$ adalah surjeksi
dari $\Cal P\Bbb N$ ke $H$.\ \Qed}

\leader{2A1I}{Filter} \cmmnt{Saya berhenti sejenak untuk membahas suatu
konstruksi yang sangat berguna dalam menyelidiki ruang topologis, tetapi juga
mempunyai kegunaan lain, dan menurut sifatnya termasuk dalam teori himpunan
elementer (bahkan jauh lebih elementer daripada pekerjaan di atas).

\header{2A1Ia}}{\bf (a)} Misalkan $X$ suatu himpunan tak kosong. Suatu
{\bf filter} pada $X$ adalah keluarga $\Cal F$ dari subhimpunan-subhimpunan
$X$ sedemikian sehingga

\inset{$X\in\Cal F$,\quad $\emptyset\notin\Cal F$,

$E\cap F\in\Cal F$ kapan pun $E$, $F\in\Cal F$,

$E\in\Cal F$ kapan pun $X\supseteq E\supseteq F\in\Cal F$.}

\noindent Syarat kedua mengakibatkan (dengan induksi pada $n$) bahwa
$F_0\cap\ldots\cap F_n\in\Cal F$ kapan pun
$F_0,\ldots,F_n\in\Cal F$.

\header{2A1Ib}{\bf (b)} Misalkan $X$, $Y$ himpunan tak kosong, $\Cal F$
suatu filter pada $X$ dan $f:D\to Y$ suatu fungsi, dengan $D\in\Cal F$.
Maka

\Centerline{$\{E:E\subseteq Y,\,f^{-1}[E]\in\Cal F\}$}

\noindent adalah suatu filter pada $Y$\prooflet{ (sebab $f^{-1}[Y]=D$,
$f^{-1}[\emptyset]=\emptyset$,
$f^{-1}[E\cap F]=f^{-1}[E]\cap f^{-1}[F]$,
$X\supseteq f^{-1}[E]\supseteq f^{-1}[F]$ kapan pun
$Y\supseteq E\supseteq F$)}; saya akan menyebutnya $f[[\Cal F]]$,
{\bf filter citra} dari $\Cal F$ di bawah $f$.

\cmmnt{\medskip

\noindent{\bf Catatan} Tentu saja terdapat suatu variabel tersembunyi dalam
notasi ini. Biasanya dalam buku ini saya menganggap suatu fungsi $f$
didefinisikan oleh domainnya $\dom f$ dan nilai-nilainya pada domain itu;
yakni, fungsi tersebut ditentukan oleh grafiknya
$\{(x,f(x)):x\in\dom f\}$, dan memang biasanya saya tidak membedakan suatu
fungsi dari grafiknya. Ini berarti bahwa ketika saya menuliskan
`$f:D\to Y$ adalah suatu fungsi', kelas $D=\dom f$ dapat dipulihkan dari
fungsi tersebut, tetapi kelas $Y$ tidak; yang saya janjikan hanyalah bahwa
$Y$ memuat kelas nilai $f[D]$ dari $f$. Sekarang, dalam notasi
$f[[\Cal F]]$ di atas, kita memang perlu mengetahui himpunan $Y$ yang
menjadi tempat filter itu, walaupun hal ini tidak dapat diketahui dari
$f$ dan $\Cal F$. Jadi Anda harus selalu menyimpulkannya dari konteks.
}

\leader{2A1J}{Aksioma Pilihan}\cmmnt{ Sekarang saya sampai pada paruh kedua
bagian ini, tempat saya membahas konsep dan teorema yang bergantung pada
Aksioma Pilihan.} Izinkan saya mengingatkan Anda akan pernyataan aksioma ini:

\inset{\hskip-20pt (AC) `setiap kali $I$ merupakan suatu himpunan dan
$\familyiI{X_i}$ suatu keluarga himpunan tak kosong yang diindeks oleh $I$,
terdapat suatu fungsi $f$, dengan domain $I$, sedemikian sehingga
$f(i)\in X_i$ untuk setiap $i\in I$'.}

\noindent Fungsi $f$ adalah suatu {\bf fungsi pilihan}\cmmnt{; fungsi ini
memilih satu anggota dari setiap anggota keluarga himpunan tak kosong $X_i$
yang diberikan}.

\cmmnt{Saya percaya bahwa sikap seseorang terhadap prinsip ini merupakan
perkara pilihan pribadi. Prinsip ini adalah landasan yang mutlak diperlukan
bagi bagian-bagian sangat besar dari matematika murni abad kedua puluh,
termasuk bagian yang cukup besar dari jilid ini; tetapi terdapat pula bidang-
bidang penting tempat prinsip-prinsip yang justru bertentangan dengannya
dapat digunakan dengan hasil mencolok, dan---menurut pandangan saya---mengarah
kepada matematika yang sama sahnya. (Saya akan menguraikan salah satunya dalam
\S567 dari Jilid 5.)
Saat ini, jika diukur menurut volumenya, lebih banyak kegiatan matematika
mutakhir bergantung pada penegasan aksioma pilihan daripada pada seluruh
pesaingnya digabungkan; tetapi tempat ditemukannya gagasan paling penting
atau menggairahkan merupakan perkara penilaian dan selera. Untuk jilid ini
saya mengikuti praktik baku dalam analisis abstrak abad kedua puluh, dengan
menggunakan aksioma pilihan kapan pun diperlukan.}

\vleader{60pt}{2A1K}{Teorema Pengurutan Baik Zermelo (a)} Aksioma Pilihan
ekuivalen dengan setiap pernyataan

\inset{`untuk setiap himpunan $X$ terdapat suatu pengurutan baik dari $X$',}

\inset{`untuk setiap himpunan $X$ terdapat bijeksi antara $X$ dan suatu
ordinal',}

\inset{`untuk setiap himpunan $X$ terdapat tepat satu ordinal awal $\kappa$
sedemikian sehingga terdapat bijeksi antara $X$ dan $\kappa$.'}

\prooflet{\noindent(\Enderton, hlm.\ 196 dan seterusnya; \Halmos, \S17;
\Henle, 9.1--9.3; \Krivine, hlm.\ 20; \Lipschutz, 12.1; \Roitman, 3.6.38.)}

\header{2A1Kb}{\bf (b)} Dengan mengasumsikan aksioma pilihan,\cmmnt{
seperti yang saya lakukan hampir di seluruh risalah ini,} saya menuliskan
$\#(X)$ untuk ordinal awal $\kappa$ sedemikian sehingga terdapat bijeksi
antara $\kappa$ dan $X$; saya menyebutnya {\bf kardinal} dari $X$.

\leader{2A1L}{Akibat-akibat mendasar Aksioma Pilihan (a)} Untuk setiap dua
himpunan $X$ dan $Y$, terdapat bijeksi antara $X$ dan $Y$ jika dan hanya jika
$\#(X)=\#(Y)$. Secara lebih umum, terdapat injeksi dari $X$ ke $Y$ jika dan
hanya jika $\#(X)\le\#(Y)$, dan terdapat surjeksi dari $X$ ke seluruh $Y$
jika dan hanya jika $\#(X)\ge\#(Y)>0$ atau $\#(X)=\#(Y)=0$.

\spheader 2A1Lb\cmmnt{ Khususnya,} $\#(\Cal P\Bbb N)=\#(\Bbb R)$;
tuliskan $\frak c$ untuk nilai bersama ini, yakni {\bf kardinal kontinum}.
\cmmnt{Teorema Cantor bahwa $\Cal P\Bbb N$ dan $\Bbb R$ tak terhitung
menjadi hasil $\omega<\frak c$, yaitu,} $\omega_1\le\frak c$.

\header{2A1Lc}{\bf (c)} Jika $X$ sebarang himpunan tak hingga, dan $r\ge 1$,
maka terdapat bijeksi antara $X^r$ dan $X$. \prooflet{(\Enderton, hlm.\ 162;
\Halmos, \S24.)} \cmmnt{(Saya mencatat bahwa kita memerlukan suatu bentuk
aksioma pilihan untuk membuktikan hasil ini pada tingkat keumuman tersebut.
Namun tentu saja, untuk kebanyakan himpunan tak hingga yang muncul secara
alami dalam matematika---himpunan seperti $\Bbb N$ dan $\Cal P\Bbb R$---
hasilnya mudah dibuktikan tanpa mengacu pada aksioma pilihan.)}

\header{2A1Ld}{\bf (d)} Andaikan $\kappa$ suatu kardinal tak hingga. Jika
$I$ adalah himpunan dengan kardinal paling besar $\kappa$ dan
$\langle A_i\rangle_{i\in I}$ merupakan suatu keluarga himpunan dengan
$\#(A_i)\le\kappa$ untuk setiap $i\in I$, maka
$\#(\bigcup_{i\in I}A_i)\le\kappa$. Akibatnya
$\#(\bigcup\Cal A)\le\kappa$ kapan pun $\Cal A$ adalah keluarga himpunan
sedemikian sehingga $\#(\Cal A)\le\kappa$ dan $\#(A)\le\kappa$ untuk setiap
$A\in\Cal A$. \cmmnt{Khususnya,} $\omega_1$ tidak dapat dinyatakan sebagai
gabungan terhitung himpunan-himpunan terhitung, dan $\omega_2$ tidak dapat
dinyatakan sebagai gabungan terhitung himpunan-himpunan dengan kardinal
paling besar $\omega_1$.

\header{2A1Le}{\bf (e)} Sekarang kita dapat menyatakan ulang 2A1Hc sebagai:
jika $\#(X)\le\frak c$, maka $\#(H)\le\frak c$, dengan $H$ adalah himpunan
fungsi dari suatu subhimpunan terhitung dari $X$ ke $\{0,1\}$.
\prooflet{\Prf\ Sebab kita mempunyai injeksi dari $X$ ke $\Cal P\Bbb N$,
dan karena itu suatu surjeksi dari $\Cal P\Bbb N$ ke $H$.\ \Qed}

\spheader 2A1Lf Setiap kelas tak kosong dari kardinal mempunyai anggota
terkecil\cmmnt{ (menurut 2A1Dc)}.

\leader{2A1M}{Lema Zorn}\cmmnt{ Dalam 2A1K saya menguraikan prinsip
pengurutan baik.} Sekarang saya sampai pada proposisi lain yang ekuivalen
dengan aksioma pilihan:

\inset{`Misalkan $(P,\le)$ suatu himpunan terurut parsial tak kosong
sedemikian sehingga setiap subhimpunan terurut total tak kosong dari $P$
mempunyai batas atas dalam $P$. Maka $P$ mempunyai suatu elemen maksimal.'}

\noindent Inilah {\bf Lema Zorn}. \prooflet{Untuk bukti bahwa aksioma pilihan
mengakibatkan, dan diakibatkan oleh, Lema Zorn, lihat \Enderton, hlm.\ 151;
\Halmos, \S16; \Henle, 9.1--9.3; \Roitman, 3.6.38.}

\leader{2A1N}{Ultrafilter} Suatu filter $\Cal F$ pada suatu himpunan $X$
adalah {\bf ultrafilter} jika untuk setiap $A\subseteq X$ berlaku
$A\in\Cal F$ atau $X\setminus A\in\Cal F$.

Jika $\Cal F$ adalah ultrafilter pada $X$ dan $f:D\to Y$ suatu fungsi,
dengan $D\in\Cal F$, maka $f[[\Cal F]]$ adalah ultrafilter pada
$Y$\prooflet{ (sebab $f^{-1}[Y\setminus A]=D\setminus f^{-1}[A]$ untuk
setiap $A\subseteq Y$)}.

Satu jenis ultrafilter dapat diuraikan dengan mudah: jika $x$ adalah
sebarang titik dari suatu himpunan $X$, maka
$\Cal F=\{F:x\in F\subseteq X\}$ adalah ultrafilter pada $X$.
(\prooflet{Anda hanya perlu membaca definisi-definisinya. }Ultrafilter jenis
ini disebut {\bf ultrafilter prinsipal}.)
\cmmnt{Tetapi tidaklah jelas bahwa ada ultrafilter lain, dan memang tidak
mungkin membuktikan keberadaannya tanpa menggunakan suatu bentuk kuat dari
aksioma pilihan, sebagai berikut.}

\leader{2A1O}{Teorema Ultrafilter}\cmmnt{ Sebagai suatu contoh penggunaan
Lema Zorn yang akan sangat berharga dalam mempelajari ruang topologis kompak
(2A3N {\it dan seterusnya}, serta \S247), saya memberikan hasil berikut.

\medskip

\noindent{\bf Teorema}} Misalkan $X$ sebarang himpunan tak kosong, dan
$\Cal F$ suatu filter pada $X$. Maka terdapat suatu ultrafilter $\Cal H$
pada $X$ sedemikian sehingga $\Cal F\subseteq\Cal H$.

\proof{ (Bandingkan \Henle, 9.4; \Roitman, 3.6.37.) Misalkan $\frak P$
adalah himpunan semua filter pada $X$ yang memuat $\Cal F$, dan urutkan
$\frak P$ menurut inklusi, sehingga untuk $\Cal G_1$,
$\Cal G_2\in\frak P$, berlaku $\Cal G_1\le \Cal G_2$ dalam $\frak P$ jika
dan hanya jika $\Cal G_1\subseteq\Cal G_2$. Mudah dilihat bahwa $\frak P$
merupakan himpunan terurut parsial, dan himpunan itu tak kosong karena
$\Cal F\in\frak P$. Jika $\frak Q$ sebarang subhimpunan terurut total tak
kosong dari $\frak P$, maka
$\Cal H_{\frak Q}=\bigcup\frak Q\in\frak P$. \Prf\ Tentu saja
$\Cal H_{\frak Q}$ merupakan suatu keluarga subhimpunan $X$. (i) Ambil
sebarang $\Cal G_0\in\frak Q$; maka
$X\in\Cal G_0\subseteq\Cal H_{\frak Q}$. Jika $\Cal G\in\frak Q$, maka
$\Cal G$ suatu filter, sehingga $\emptyset\notin\Cal G$; karena itu
$\emptyset\notin\Cal H_{\frak Q}$. (ii) Jika $E$,
$F\in\Cal H_{\frak Q}$, maka terdapat $\Cal G_1$, $\Cal G_2\in\frak Q$
sedemikian sehingga $E\in\Cal G_1$ dan $F\in\Cal G_2$. Karena $\frak Q$
terurut total, berlaku $\Cal G_1\subseteq\Cal G_2$ atau
$\Cal G_2\subseteq\Cal G_1$. Dalam kedua kasus,
$\Cal G=\Cal G_1\cup\Cal G_2\in\frak Q$. Sekarang $\Cal G$ adalah suatu
filter yang memuat $E$ dan $F$, sehingga memuat $E\cap F$, dan
$E\cap F\in\Cal H_{\frak Q}$. (iii) Jika
$X\supseteq E\supseteq F\in\Cal H_{\frak Q}$, terdapat
$\Cal G\in\frak Q$ sedemikian sehingga $F\in\Cal G$; dan
$E\in\Cal G\subseteq\Cal H_{\frak Q}$. Ini menunjukkan bahwa
$\Cal H_{\frak Q}$ adalah suatu filter pada $X$. (iv) Akhirnya,
$\Cal H_{\frak Q}\supseteq\Cal G_0\supseteq\Cal F$, sehingga
$\Cal H_{\frak Q}\in\frak P$.\ \QeD\ Sekarang $\Cal H_{\frak Q}$ jelas
merupakan batas atas bagi $\frak Q$ dalam $\frak P$.

Karena itu kita dapat menerapkan Lema Zorn untuk menemukan suatu elemen
maksimal $\Cal H$ dari $\frak P$. $\Cal H$ ini tentu merupakan filter pada
$X$ yang memuat $\Cal F$.

Sekarang misalkan $A\subseteq X$ sedemikian sehingga $A\notin\Cal H$.
Pertimbangkan

\Centerline{$\Cal H_1=\{E:E\subseteq X,\,E\cup A\in\Cal H\}$.}

\noindent Ini merupakan suatu filter pada $X$. \Prf\ Tentu saja ini adalah
suatu keluarga subhimpunan $X$. (i) $X\cup A=X\in\Cal H$, sehingga
$X\in\Cal H_1$. $\emptyset\cup A=A\notin\Cal H$, sehingga
$\emptyset\notin\Cal H_1$. (ii) Jika $E$, $F\in\Cal H_1$, maka

\Centerline{$(E\cap F)\cup A=(E\cup A)\cap(F\cup A)
\in\Cal H$,}

\noindent sehingga $E\cap F\in\Cal H_1$. (iii) Jika
$X\supseteq E\supseteq F\in\Cal H_1$, maka
$E\cup A\supseteq F\cup A\in\Cal H$, sehingga $E\cup A\in\Cal H$ dan
$E\in\Cal H_1$.\ \QeD\ Juga $\Cal H_1\supseteq\Cal H$, sehingga
$\Cal H_1\in\frak P$. Tetapi $\Cal H$ merupakan elemen maksimal dari
$\frak P$, sehingga $\Cal H_1=\Cal H$. Karena
$(X\setminus A)\cup A=X\in\Cal H$, maka $X\setminus A\in\Cal H_1$ dan
$X\setminus A\in\Cal H$.

Karena $A$ sebarang, $\Cal H$ adalah suatu ultrafilter, sebagaimana diminta.
}%end of proof of 2A1O



\leader{2A1P}{}\cmmnt{ Sekarang saya sampai pada suatu hasil dari
kombinatorika infinitari yang saya buktikan secara terperinci, bukan karena
hasil ini tidak dapat ditemukan dalam banyak buku ajar, melainkan karena
biasanya hasil tersebut diberikan pada tingkat keumuman yang jauh lebih
besar, sampai-sampai memahami mengapa teorema yang dinyatakan mencakup hasil
sekarang mungkin lebih sulit daripada membuktikan hasil terakhir ini dari
prinsip-prinsip pertama.

\medskip

\noindent}{\bf Teorema} (a) Misalkan
$\langle K_{\alpha}\rangle_{\alpha\in A}$ suatu keluarga himpunan terhitung,
dengan $\#(A)$ lebih besar daripada $\frak c$, kardinal kontinum.
Maka terdapat suatu himpunan $M$, dengan kardinal paling besar $\frak c$,
dan suatu himpunan $B\subseteq A$, dengan kardinal lebih besar
daripada $\frak c$, sedemikian sehingga
$K_{\alpha}\cap K_{\beta}\subseteq M$ kapan pun $\alpha$, $\beta$ adalah
anggota-anggota berbeda dari $B$.

(b) Misalkan $I$ suatu himpunan, dan
$\langle f_{\alpha}\rangle_{\alpha\in A}$ suatu keluarga dalam
$\{0,1\}^I$, yaitu himpunan fungsi dari $I$ ke $\{0,1\}$, dengan
$\#(A)>\frak c$. Jika $\langle K_{\alpha}\rangle_{\alpha\in A}$ sebarang
keluarga subhimpunan terhitung dari $I$, maka terdapat suatu himpunan
$B\subseteq A$, dengan kardinal lebih besar daripada $\frak c$, sedemikian
sehingga $f_{\alpha}$ dan $f_{\beta}$ sama pada
$K_{\alpha}\cap K_{\beta}$ untuk semua $\alpha$, $\beta\in B$.

(c) Khususnya, di bawah syarat-syarat (b), terdapat $\alpha$,
$\beta\in A$ yang berbeda sedemikian sehingga $f_{\alpha}$ dan $f_{\beta}$
sama pada $K_{\alpha}\cap K_{\beta}$.

\proof{{\bf (a)} Pilih secara induktif suatu keluarga
$\langle M_{\xi}\rangle_{\xi<\omega_1}$ dari himpunan-himpunan menurut aturan

\inset{jika terdapat suatu himpunan $N$ sedemikian sehingga

\inset{($*$) $N$ saling lepas dari $\bigcup_{\eta<\xi}M_{\eta}$,
$\#(N)\le\frak c$ dan
%
%\qquad\qquad
$\#(\{\alpha:\alpha\in A,\,K_{\alpha}\cap N=\emptyset\})
  \le\frak c$,}

\noindent pilihlah himpunan seperti itu sebagai $M_{\xi}$;

jika tidak, tetapkan $M_{\xi}=\emptyset$.}

\noindent Setelah $M_{\xi}$ dipilih untuk setiap $\xi<\omega_1$, tetapkan
$M=\bigcup_{\xi<\omega_1}M_{\xi}$. Aturan tersebut memastikan bahwa
$\ofamily{\xi}{\omega_1}{M_{\xi}}$ saling lepas dan bahwa
$\#(M_{\xi})\le\frak c$ untuk setiap $\xi<\omega_1$, sedangkan
$\omega_1\le\frak c$, sehingga $\#(M)\le\frak c$.

Misalkan $\frak P$ adalah keluarga himpunan $P\subseteq A$ sedemikian sehingga
$K_{\alpha}\cap K_{\beta}\subseteq M$ untuk semua $\alpha$, $\beta\in P$
yang berbeda. Urutkan $\frak P$ menurut inklusi, sehingga himpunan itu
merupakan himpunan terurut parsial. Jika $\frak Q\subseteq\frak P$ terurut
total, maka $\bigcup\frak Q\in\frak P$.
\Prf\ Jika $\alpha$, $\beta$ adalah anggota berbeda dari $\bigcup\frak Q$,
terdapat $Q_1$, $Q_2\in\frak Q$ sedemikian sehingga $\alpha\in Q_1$,
$\beta\in Q_2$; sekarang $P=Q_1\cup Q_2$ sama dengan salah satu dari $Q_1$,
$Q_2$, dan dalam kedua kasus termasuk dalam $\frak P$ serta memuat baik
$\alpha$ maupun $\beta$, sehingga
$K_{\alpha}\cap K_{\beta}\subseteq M$.\ \QeD\ Menurut Lema Zorn,
$\frak P$ mempunyai suatu elemen maksimal $B$, dan tentu saja kita mempunyai
$K_{\alpha}\cap K_{\beta}\subseteq M$ untuk semua $\alpha$,
$\beta\in B$ yang berbeda.

\Quer\ Andaikan, jika mungkin, bahwa $\#(B)\le\frak c$. Tetapkan
$N=\bigcup_{\alpha\in B}K_{\alpha}\setminus M$. Maka $N$ mempunyai kardinal
paling besar $\frak c$, karena termuat dalam gabungan paling banyak
$\frak c$ himpunan terhitung. Untuk setiap $\gamma\in A\setminus B$,
$B\cup\{\gamma\}\notin\frak P$, sehingga harus terdapat suatu $\alpha\in B$
sedemikian sehingga $K_{\alpha}\cap K_{\gamma}\not\subseteq M$; yaitu,
$K_{\gamma}\cap N\ne\emptyset$. Jadi
$\{\gamma:K_{\gamma}\cap N=\emptyset\}\subseteq B$ mempunyai kardinal
paling besar $\frak c$. Tetapi ini berarti bahwa dalam aturan untuk memilih
$M_{\xi}$, selalu terdapat suatu $N$ yang memenuhi syarat ($*$), dan karena
itu $M_{\xi}$ juga memenuhinya. Jadi
$C_{\xi}=\{\alpha:K_{\alpha}\cap M_{\xi}=\emptyset\}$ mempunyai kardinal
paling besar $\frak c$ untuk setiap $\xi<\omega_1$. Maka
$C=\bigcup_{\xi<\omega_1}C_{\xi}$ juga demikian. Tetapi hipotesis semula
adalah $\#(A)>\frak c$, sehingga terdapat suatu $\alpha\in A\setminus C$.
Dalam hal ini, $K_{\alpha}\cap M_{\xi}\ne\emptyset$ untuk setiap
$\xi<\omega_1$. Namun ini berarti bahwa kita mempunyai suatu surjeksi
$\phi:K_{\alpha}\cap M\to\omega_1$ yang diberikan dengan menetapkan

\Centerline{$\phi(i)=\xi$ jika $i\in K_{\alpha}\cap M_{\xi}$.}

\noindent Karena $\#(K_{\alpha})\le\omega<\omega_1$, ini mustahil.\ \Bang

Karena itu $\#(B)>\frak c$ dan kita telah menemukan pasangan $M$, $B$ yang
sesuai.

\medskip

{\bf (b)} Menurut (a), kita dapat menemukan suatu himpunan $M$, dengan
kardinal paling besar $\frak c$, dan suatu himpunan $B_0\subseteq A$, dengan
kardinal lebih besar daripada $\frak c$, sedemikian sehingga
$K_{\alpha}\cap K_{\beta}\subseteq M$ untuk semua $\alpha$,
$\beta\in B_0$ yang berbeda. Misalkan $H$ adalah himpunan fungsi dari
subhimpunan-subhimpunan terhitung dari $M$ ke $\{0,1\}$; maka
$f'_{\alpha}=f_{\alpha}\restr(K_{\alpha}\cap M)\in H$ untuk setiap
$\alpha\in B_0$. Sekarang
$B_0=\bigcup_{h\in H}\{\alpha:\alpha\in B_0,\,f'_{\alpha}=h\}$ mempunyai
kardinal lebih besar daripada $\frak c$, sedangkan $\#(H)\le\frak c$
(2A1Le), sehingga harus terdapat suatu $h\in H$ sedemikian sehingga
$B=\{\alpha:\alpha\in B_0,\,f'_{\alpha}=h\}$ mempunyai kardinal lebih besar
daripada $\frak c$.

Jika $\alpha$, $\beta$ adalah anggota berbeda dari $B$, maka
$K_{\alpha}\cap K_{\beta}\subseteq M$, sebab $\alpha$, $\beta\in B_0$;
tetapi ini berarti bahwa

\Centerline{$f_{\alpha}\restr K_{\alpha}\cap K_{\beta}
=h\restr K_{\alpha}\cap K_{\beta}
=f_{\beta}\restr K_{\alpha}\cap K_{\beta}$.}

\noindent Jadi $B$ mempunyai sifat yang diperlukan. (Tentu saja
$f_{\alpha}$ dan $f_{\beta}$ sama pada $K_{\alpha}\cap K_{\beta}$ jika
$\alpha=\beta$.)

\medskip

{\bf (c)} langsung mengikuti.
}%end of proof of 2A1P
%elementary submodel argument

\cmmnt{\medskip

\noindent{\bf Catatan} Hasil yang kita perlukan dalam jilid ini (dalam 216E)
adalah bagian (c) di atas. Ada bukti-bukti lain untuk ini, mungkin sedikit
lebih sederhana; tetapi hasil yang lebih kuat dalam bagian (b) akan berguna
dalam Jilid 3.
}%end of comment

\discrpage
